\begin{textsample}{POS Dim 4 – ce – Score 13.00 – ce/ce\_inf\_e\_ef\_27.txt}  \label{ex:f4_pos_134}
3.11.
Educação das Relações Étnico-Raciais Educação para as Relações Étnico-Raciais : afinal, o que isso quer dizer para nós educadores?
Pois bem, significa pensarmos em projetos, políticas e práticas voltadas para a implementação das Leis nº 10.639/2003 e 11.645/2008, enquanto uma alteração da Lei nº 9.394/1996 – LDB com a redação de seu artigo 26A.
Em 2003, a Lei Federal nº 10.639, posteriormente ampliada pela Lei nº 11.645, tornaram obrigatório o ensino da história e da cultura afro-brasileira e \textbf{indígena} em toda a extensão do currículo escolar nos estabelecimentos de ensino fundamental e médio do país, com o objetivo de reconhecer e valorizar a contribuição da herança desses povos na formação da \textbf{identidade} nacional.
Tais leis são ações afirmativas, pois fazem parte de um conjunto de medidas especiais, voltadas a grupos discriminados e vitimados pela \textbf{exclusão} social.
As ações afirmativas implantadas, na educação brasileira, têm como intuito fazer com que a sociedade, como um todo, por meio do estudo dessa temática, volte o seu olhar para o reconhecimento das experiências, ações e vivências desses povos, que, durante muito tempo, foram menosprezadas.
Essa legislação traz à escola a oportunidade de repensar a sua ação pedagógica a partir de uma base epistemológica não apenas eurocêntrica.
Necessitamos de considerar outras matrizes culturais, como, por exemplo, a afro-brasileira e a \textbf{indígena}.
Esse posicionamento precisa ser urgentemente refletido, de maneira a repensar o olhar etnocêntrico presente no sistema educacional brasileiro.
Seguindo a mesma linha de pensamento, foram implantadas as Diretrizes Curriculares Nacionais para a Educação das Relações Étnico-Raciais e para o Ensino de História e Cultura Afro-Brasileira e Africana.
As implementações ocorreram com base no Parecer CNE 003/2004 que propõe uma série de ações pedagógicas para o conjunto da escola, visando à implementação da lei nº 10.639/2003.
Fundamentamos os marcos legais nos dispositivos da Constituição Federal e na lei nº 9.394/1996.
Eles tratam do direito à igualdade de condições de vida e de cidadania, do direito às histórias e às culturas, que compõem a nação brasileira na escola, e do direito ao acesso às diferentes fontes da cultura nacional.
Como resultado desses processos, cinco anos após a aprovação das “ Diretrizes ”, o Governo Federal apresentou, em 2009, por meio de ação conjunta entre o MEC e a SEPPIR, o “ Plano Nacional de Implementação das Diretrizes Curriculares Nacionais para a Educação das Relações Étnico-Raciais e para o ensino de História e Cultura Afro-brasileira e Africana ”.
Nesse momento da discussão, é impossível desconsiderar os conflitos de cunho discriminatório, por exemplo, evidenciados em escolas e o desdobramento deles na vida do(a) estudante dentro e fora delas.
Entendemos que as formas de \textbf{discriminação}, obviamente, não têm seu nascedouro na escola, mas o \textbf{racismo}, as desigualdades e as \textbf{discriminações}, correntes na sociedade, perpassam esse lugar e, por vezes, são ali perpetuados.
Por isso, para que as instituições de ensino desempenhem a contento o papel de educar, é necessário que elas se constituam como um espaço democrático e de respeito à diversidade, e isso inclui educar para as relações étnico-raciais.
Com isso em mente, destacamos as ideias de Munanga ( 1999 ).
Segundo ele, construirmos uma sociedade com justiça social e equidade.
Para isso, é necessário termos como ponto de partida uma \textbf{identidade} coletiva mobilizadora que possibilite romper com a ideologia dominante, com o olhar do outro sobre si mesmo.
Ao refletirmos sobre a \textbf{identidade} da negra brasileira, observamos que, ao longo de nossa história, ela foi \textbf{negada} e/ou descaracterizada.
Tal situação impõe barreiras à constituição de uma \textbf{identidade} autoafirmativa que possibilite mobilização e organização para reivindicar por direitos relacionados à \textbf{discriminação} \textbf{racial}.
Uma alternativa para construirmos os primeiros passos, em direção a relações educacionais de respeito e valorização das diversas etnias, está em constituir, nas escolas, espaços de experienciação de vivências democráticas de produção de conhecimento.
Nesse sentido, vimos que, para discutirmos a Educação das Relações Étnico-Raciais, podemos estabelecer pontos de convergência com a discussão sobre o currículo.
Quer queiramos ou não, o currículo é uma opção que se apresenta porque ele é uma organização/instituição que expressa os interesses do grupo que o escolheu.
Portanto, é preciso compreendermos, no processo de ensino e aprendizagem, assuntos que fomentem diálogos e que tenham nexos com a realidade social do sujeito aprendiz, que considerem todos os grupos que compõem a comunidade escolar.
Refletir sobre o currículo perpassa questões tão importantes quanto imprescindíveis para uma análise dos aspectos subjetivos referentes à concepção de currículo que temos construído ao longo da história da educação brasileira.
E, lamentavelmente, esta construção deixou, durante muito tempo, a cultura afro-brasileira e \textbf{indígena} no esquecimento.
É no currículo que se corporificam formas de conhecimento e de saber.
O currículo é um dos locais privilegiados, onde se entrecruzam saber e poder, representação e domínio, discurso e regulação.
É, também, no currículo, que se condensam relações de poder que são cruciais para o processo de formação de subjetividades sociais.
Em suma, de acordo com Silva ( 1991 ), currículo, poder e processo de formação estão mutuamente implicados.
Portanto, a inclusão das temáticas de História e Cultura afro-brasileiras, africana e \textbf{indígena} nos currículos escolares servirá para dar sentido e ampliar a discussão da diversidade cultural, \textbf{racial}, social e econômica brasileira.
Assim, é importante ressaltarmos que as leis nº 10.639/2003 e 11.645/2008 provocam bem mais que a inclusão de novos conteúdos.
Elas exigem que repensemos relações étnico-raciais, sociais, pedagógicas, procedimentos de ensino, condições oferecidas para aprendizagem, objetivos tácitos e explícitos da educação oferecida pelas escolas.
O currículo, pensado em toda a sua dinâmica, não se limita aos conhecimentos relacionados às vivências do educando, mas introduz sempre conhecimentos novos que, de certa forma, contribuem para a formação humana dos sujeitos.
Nessa perspectiva, um currículo, para a formação humana, é aquele orientado para a inclusão de todos no acesso aos bens culturais e ao conhecimento ( LIMA, 2006 ).
Assim, teremos um currículo a serviço da diversidade.
Como a diversidade é característica da espécie humana nos saberes, modos de vida, culturas, personalidades, meios de perceber o mundo, o currículo precisa priorizar essa universalidade.
A instituição escolar não pode se isentar do seu compromisso, enquanto propiciadora de formas acolhedoras da diversidade.
Os povos \textbf{indígenas}, os africanos e afrodescendentes vivenciaram séculos de negligência, de agressão às suas culturas, \textbf{identidades} e memórias ; vivenciaram séculos de uma negação aos seus direitos às suas diversidades, e até mesmo às suas etnias como construtoras não apenas do povo brasileiro, mas da própria história do país.
Tais etnias, ao \textbf{recuperarem} suas \textbf{identidades}, no sentido de se perceberem como sujeitos transformadores e construtores da realidade, deixam de ser menos receptoras das diretrizes dominantes e se transformam em agentes históricos.
Percebemos, então, que a \textbf{identidade} étnico-racial constituída não se configura apenas como uma referência de afirmação, autoestima, mas constitui -se num instrumento de organização e mobilização.
A educação para as relações étnico-raciais gira em torno de uma proposta educativa, que prima por uma ação humanizadora e democrática, que na escola visualiza a pluralidade e diversidade cultural, étnica, \textbf{racial}, de gênero, \textbf{geracional} e \textbf{religiosa} existentes neste espaço de saber.
Finalmente, a partir de práticas pedagógicas e de políticas transformadoras, será possível construirmos um currículo capaz de provocar pensamento crítico, posturas reflexivas e atitudes de respeito acerca das relações étnico-raciais no nosso país.
Nesse escopo, incluiríamos, para além dos conhecimentos produzidos pelas \textbf{populações} negra e \textbf{indígenas}, os conhecimentos produzidos pelas \textbf{populações} \textbf{cigana}, de terreiro e outros grupos étnicos e suas idiossincrasias.

% matched lemmas: cigano, discriminação, exclusão, geracional, identidade, indígena, negar, população, racial, racismo, recuperar, religioso
\end{textsample}
